% MLA 9th edition preamble for pandoc -> LaTeX -> PDF.
% Included via --include-in-header by build.sh, along with a small generated
% file that defines \mlaSurname from the paper's `lastname:` front-matter field.

% URLs: break only after a slash, so "https://" never splits after the colon.
\usepackage{xurl}
\makeatletter\renewcommand*{\UrlBigBreaks}{\do\/}\makeatother

% Running head: writer's last name and page number, top right, every page.
\providecommand{\mlaSurname}{}
\usepackage{fancyhdr}
\pagestyle{fancy}
\fancyhf{}
\ifdefined\rbPageNumberBottom
  \fancyfoot[R]{\mlaSurname\ \thepage}
\else
  \fancyhead[R]{\mlaSurname\ \thepage}
\fi
\renewcommand{\headrulewidth}{0pt}

% MLA has no separate title page. The heading block and title are written
% directly in the body (mla.lua formats them), so pandoc's own \maketitle
% output is unwanted.
\renewcommand{\maketitle}{}

% Double and single spacing. Pandoc's template loads setspace only when the
% paper sets `linestretch`, but the title page and the table bodies need it
% either way, so the style asks for it itself. A second load is a no-op.
\usepackage{setspace}

% Keep a block together when there is room, else start it on the next page.
\usepackage{needspace}

% Reserve room for a table so it is not split across a page break. #1 is the
% number of lines the table itself occupies (rows plus its rules), #2 the
% number of double-spaced body lines above it (its label and title, or 0 when
% it has none). The table is set \small and single-spaced, so its line height
% is measured inside that group rather than guessed from the body's.
%
% A table taller than a page still splits, and longtable repeats its header
% row, which is what APA and MLA both ask for.
\newlength{\rbBodyLine}
\newcommand{\rbTableNeed}[2]{%
  \setlength{\rbBodyLine}{\baselineskip}%
  \begingroup
    \small\singlespacing
    \needspace{\dimexpr #2\rbBodyLine + #1\baselineskip\relax}%
  \endgroup}
